% ============================================================================
% CHAPTER 10 — CROSS-APPROACH CONNECTIONS AND FIELD TRAJECTORY
% ============================================================================
\chapter{Cross-Approach Connections and Future Directions}
\label{ch:future}

\begin{motivation}
    Having compared the individual approaches, we now step back further
    to examine how these methods relate to each other historically and
    conceptually, what recurring themes emerge, what problems remain
    unsolved, and where the field is heading.
\end{motivation}

\section{How the Approaches Build on Each Other}
\label{sec:evolution}

The six approaches studied in this document did not emerge independently---each
builds on insights from earlier work.
Figure~\ref{fig:evolution} illustrates both the chronological development and
the conceptual dependency structure.

\begin{figure}[H]
    \centering
    \begin{tikzpicture}[
        % --- Node & Arrow styles ---
        base/.style={
            rounded corners=4pt,
            font=\footnotesize\bfseries,
            minimum height=1.8em,
            minimum width=5.5em,
            inner sep=4pt,
            text=white,
            align=center
        },
        slm/.style={base, fill=navyDark},
        ae/.style={base, fill=motivationColor},
        scf/.style={base, fill=insightColor},
        unfold/.style={base, fill=misconceptionColor},
        dep/.style={->, >=stealth, thick, color=rulegray, rounded corners=4pt},
        yearlbl/.style={font=\scriptsize\itshape, text=chaptergray}
    ]
        % --- Timeline nodes ---
        \node[slm]    (ga)     at (0, 0)     {GA-SLM};
        \node[ae]     (pr)     at (2.8, 0)   {PRNet};
        \node[ae]     (eslm)   at (5.6, 0)   {ESLM-AE};
        \node[scf]    (nn)     at (8.4, 0)   {NN-SCF};
        \node[unfold] (prdun)  at (11.2, 0)  {PR-DUN};
        \node[ae]     (dlae)   at (14.0, 0)  {DL-AE};

        % --- Year labels ---
        \foreach \name/\year in {ga/2016, pr/2017, eslm/2019, nn/2021, prdun/2022, dlae/2025} {
            \node[yearlbl] at (\name.south) [yshift=-12pt] {\year};
        }

        % --- Dependency arrows (Orthogonal straight lines, all above) ---
        % Staggered entry/exit points using angles to prevent arrow collision
        
        % 1. GA-SLM -> ESLM-AE
        \draw[dep] (ga.north)   -- ++(0, 0.4) -| (eslm.110);
        
        % 2. PRNet -> ESLM-AE (Direct horizontal, shifted up slightly)
        \draw[dep] (pr.15) -- (eslm.165);
        
        % 3. PRNet -> NN-SCF
        \draw[dep] (pr.70)   -- ++(0, 0.8) -| (nn.north);
        
        % 4. ESLM-AE -> DL-AE
        \draw[dep] (eslm.70) -- ++(0, 1.2) -| (dlae.115);
        
        % 5. PRNet -> DL-AE
        \draw[dep] (pr.110)  -- ++(0, 1.6) -| (dlae.65);

        % --- Legend ---
        \matrix[anchor=center, draw=rulegray, rounded corners=3pt,
                inner sep=5pt, column sep=6pt, row sep=2pt,
                font=\scriptsize, fill=white] at (7.0, -1.8) {
            \node[slm,    minimum width=1em, minimum height=0.8em] {}; & \node {SLM-based}; &
            \node[ae,     minimum width=1em, minimum height=0.8em] {}; & \node {Autoencoder}; &
            \node[scf,    minimum width=1em, minimum height=0.8em] {}; & \node {Signal Cancellation}; &
            \node[unfold, minimum width=1em, minimum height=0.8em] {}; & \node {Deep Unfolding}; \\
        };
    \end{tikzpicture}
    \caption{Chronological evolution and dependency structure of ML-enhanced
             PAPR reduction approaches studied in this document.
             Arrows indicate conceptual dependency (a later approach built on
             an earlier one).  Node colors indicate the approach category.}
    \label{fig:evolution}
\end{figure}

\begin{keyinsight}
    The field has evolved from \emph{augmenting} classical SLM
    (GA-SLM, ESLM-AE) toward \emph{replacing} it with learned
    alternatives (PRNet, NN-SCF, DL-AE), and most recently
    toward \emph{hybrid model-driven} designs that combine the best of
    both worlds (PR-DUN).
    This mirrors a broader trend in signal processing: from
    model-free ML $\to$ model-aware ML $\to$ model-driven
    ML~\cite{da2024survey}.
\end{keyinsight}

\section{Recurring Themes}
\label{sec:themes}

Several themes appear consistently across all approaches:

\paragraph{Theme 1: The PAPR--BER Trade-off Is Fundamental.}
Every approach must balance reducing PAPR against maintaining signal fidelity.
This manifests as the $\alpha$ parameter in PRNet and NN-SCF, the power budget $\epsilon$ in PR-DUN, and the $\beta$ parameter in ESLM-AE and DL-AE.
No approach eliminates the trade-off; they merely provide different mechanisms to control it~\cite{kim2017novel,zou2021novel,nguyen2022deep,hao2019extended,alnaseri2025deep}.

\paragraph{Theme 2: Training Data Is Free.}
All DNN-based approaches benefit enormously from the fact that OFDM symbols can be generated synthetically.
Unlike image recognition or natural language processing, the training data distribution is \emph{known exactly} and can be sampled at will. This removes a major bottleneck of deep learning~\cite{da2024survey}.

\paragraph{Theme 3: Channel Modeling Is the Weak Link.}
Most approaches (PRNet, NN-SCF, PR-DUN) evaluate under AWGN only. The few that test more realistic channels (ESLM-AE under VLC channels, DL-AE under fiber) show promising but unvalidated performance. Robust operation under diverse, time-varying channels remains an open challenge~\cite{da2024survey}.

\paragraph{Theme 4: Scalability Is Untested.}
All papers use moderate subcarrier counts ($N = 64$--$256$). Practical OFDM systems use $N = 2048$ (LTE) or $N = 4096$ (5G NR). Whether FC-layer-based approaches scale to these sizes is unclear; the number of weights grows as $\mathcal{O}(N^2)$, which may be prohibitive~\cite{da2024survey}.

\section{What Has ML Genuinely Solved?}
\label{sec:solved}

Based on the evidence from our seven source papers, ML has genuinely solved or significantly advanced the following aspects of PAPR reduction:

\begin{enumerate}
    \item \textbf{Side information elimination:}
          Autoencoders (PRNet, ESLM-AE, DL-AE) conclusively demonstrate
          that joint encoder--decoder training can eliminate SI,
          a longstanding SLM problem~\cite{kim2017novel,hao2019extended,alnaseri2025deep}.
    \item \textbf{Fixed-complexity inference:}
          DNN-based approaches replace the $U$-dependent SLM complexity
          with a single fixed-cost forward pass, making the inference
          cost independent of the desired PAPR
          reduction~\cite{kim2017novel,zou2021novel}.
    \item \textbf{Continuous phase optimization:}
          ESLM-AE showed that extending the phase alphabet to
          $[0, 2\pi)$ and learning the optimal phases via
          backpropagation outperforms discrete-alphabet SLM~\cite{hao2019extended}.
    \item \textbf{Algorithm parameter optimization:}
          PR-DUN demonstrated that even a classical iterative algorithm
          can be significantly improved by learning its
          per-iteration parameters~\cite{nguyen2022deep}.
\end{enumerate}

\section{What Problems Remain Open?}
\label{sec:open_problems}

\begin{enumerate}
    \item \textbf{Scalability to large $N$:}
          No ML approach has been demonstrated for $N > 256$.
          FC-layer architectures require $\mathcal{O}(N^2)$ parameters,
          which is impractical for $N = 2048$+.
          Convolutional or transformer-based architectures may be
          needed~\cite{da2024survey}.
    \item \textbf{Channel robustness:}
          Most approaches train and test under a single channel model.
          How do they perform under unseen channel conditions?
          Transfer learning across channel types is unexplored.
    \item \textbf{Multi-user OFDM:}
          All approaches consider single-user OFDM.
          Extension to OFDMA (multi-user) introduces inter-user
          coordination challenges.
    \item \textbf{MIMO-OFDM:}
          Modern systems use multiple antennas.
          PAPR reduction in MIMO-OFDM is a multi-dimensional
          optimization problem that none of these approaches address.
    \item \textbf{Real-time hardware implementation:}
          No paper provides FPGA or ASIC implementation results.
          The real-world feasibility of DNN inference at symbol rates
          of 28,000+ symbols/s remains unproven~\cite{da2024survey}.
    \item \textbf{Standardization and deployment:}
          None of the ML approaches have been incorporated into
          any wireless or optical communication standard.
    \item \textbf{Joint optimization with other subsystems:}
          PAPR reduction does not exist in isolation.
          Joint optimization with channel coding, equalization,
          and pre-distortion (as hinted by NN-SCF's DPD experiment)
          is a promising but largely unexplored
          direction~\cite{zou2021novel,da2024survey}.
\end{enumerate}

\section{What Would a ``Perfect'' SLM System Look Like?}
\label{sec:perfect_system}

Based on the lessons learned from all seven papers, a hypothetical ``perfect'' ML-enhanced SLM system would:

\begin{enumerate}
    \item Be \textbf{distortion-free}: phase rotation only, no additive
          signals or clipping.
    \item \textbf{Eliminate side information}: the receiver determines the
          phase sequence automatically.
    \item Have \textbf{fixed, low-complexity inference}: a single forward
          pass independent of $U$ or $N$.
    \item \textbf{Scale} to $N = 2048$+ without prohibitive parameter counts.
    \item Operate robustly across \textbf{diverse channels} without retraining.
    \item Be \textbf{implementable in hardware} at real-time symbol rates.
    \item Achieve \textbf{5+ dB PAPR reduction} at $\mathrm{CCDF} = 10^{-3}$.
\end{enumerate}

No existing approach satisfies all seven criteria. ESLM-AE comes closest (1, 2, 3), but lacks scalability (4), channel robustness (5), and hardware validation (6). Achieving all seven simultaneously is the grand challenge for future research~\cite{da2024survey}.

\section{Suggested Future Directions}
\label{sec:future_directions}

\begin{enumerate}
    \item \textbf{Lightweight architectures:}
          Replace FC layers with depthwise separable convolutions or
          attention mechanisms that scale as $\mathcal{O}(N\log N)$
          instead of $\mathcal{O}(N^2)$.
    \item \textbf{Meta-learning for channel adaptation:}
          Train a meta-model that can quickly adapt to new channel
          conditions with minimal fine-tuning data (few-shot learning).
    \item \textbf{Reinforcement learning for phase selection:}
          Frame the phase-sequence selection as a sequential decision
          problem and use RL agents to learn selection policies
          that generalize across symbols~\cite{da2024survey}.
    \item \textbf{Model-driven + data-driven hybrids:}
          Extend the PR-DUN idea to SLM-specific algorithms,
          unfolding the SLM search into a trainable network.
    \item \textbf{End-to-end OFDM learning:}
          Jointly optimize modulation, PAPR reduction, channel coding,
          and equalization in a single differentiable pipeline.
    \item \textbf{Hardware-aware training:}
          Incorporate quantization and fixed-point arithmetic constraints
          into the training process to ensure deployability on
          FPGA/ASIC platforms.
\end{enumerate}


% ----------------------------------------------------------------------------
\section{Concluding Remarks}
\label{sec:concluding_remarks}
% ----------------------------------------------------------------------------

This document has presented a comprehensive study of six ML-enhanced
approaches to the PAPR reduction problem in OFDM systems, each addressing
one or more of the three fundamental limitations of classical SLM:
computational complexity, side information overhead, and blind detection
difficulty.

The key findings can be distilled into three conclusions:

\begin{enumerate}
    \item \textbf{ML works.}
          All six approaches achieve 3--5~dB PAPR reduction, confirming that
          machine learning---whether neural networks, autoencoders, deep
          unfolding, or evolutionary optimization---can match or exceed
          classical SLM performance while offering structural advantages
          (lower complexity, no side information, or more interpretable
          operation).

    \item \textbf{No single approach solves everything.}
          Each technique makes a different trade-off between PAPR performance,
          BER impact, computational cost, and applicability domain.
          ESLM-AE comes closest to an ideal solution (distortion-free,
          SI-free, fast inference), but scalability and channel robustness
          remain open challenges for all approaches.

    \item \textbf{The future lies in hybrid, model-driven designs.}
          The field is converging toward architectures that combine
          the interpretability and structure of classical algorithms with
          the learning capability of neural networks---as exemplified by
          PR-DUN.  This hybrid philosophy is likely to produce the next
          generation of practical, deployable PAPR reduction solutions.
\end{enumerate}

As OFDM continues to underpin wireless (5G NR and beyond), optical, and
underwater communication systems, the importance of efficient PAPR reduction
will only grow.  The ML-enhanced approaches studied here represent a
promising and rapidly evolving frontier that bridges classical signal processing
with modern deep learning.
